\begin{table}

\caption{Top Earners and Losers after the Reform (in millions of Mexican pesos)}
\label{tab:earners}
\centering
\begin{tabular}[t]{cc|cc}
\hline
\textbf{Outlet (Winners)} & \textbf{Total Change} & \textbf{Outlet (Losers)} & \textbf{Total Change}\\
\hline
E y P de Medios de los Estados & 27.8 (645\%) & El Universal & -251.6 (-83\%) \\
Grupo Cantón & 25.0 (10,652\%) & El Financiero & -65.3 (-79\%) \\
El Reforma & 9.9 (41\%) & La Crónica Diaria & -63.6 (-93.5\%) \\
Proceso & 5.3 (207\%) & Expansión & -60.2 (-97\%)\\
Medios Masivos Mexicanos & 5.3 (6.9\%) & La Razón de México & -58.9 (-92\%)\\
Azteca Noticias & 4.4 (32\%) & El Sol de México & -55.5 (-91\%)\\
Editorial Golfo Pacífico & 3.8 (1979\%) & El Economista & -49.9 (-70.7\%)\\
Stereorey México & 2.8 (2172\%) & Milenio & -39.4 (-47.7\%)\\
Piedra Angular & 1.8 (55.3\%) & Información Integral & -35.5 (-87.6\%)\\
Editora de Medios Mapa & 1.7 (1250\%) & TV Notas & -34.7 (-87\%)\\
\hline \\[-1.8ex]  \multicolumn{4}{l} {\parbox[t]{18cm}{\scripsize \textit{Notes:}  Descriptive statistics of the COMSOC advertising spending dataset. The period before refers to the years 2017 and 2018, while the period after refers to 2019 and 2020. Each even column represents the total change in advertising spending from the 2017-2018 period to the 2019-2020, with the percent change it represented. All numeric variables are in millions of Mexican pesos. At that moment 1USD $\approx$ 19MXN on average.}} \\
\end{tabular}
\end{table}